\section{Details of the Spine Proofs}
\label{app:spine-proofs}

\subsection{The Bridge Invariant}
We complete the left-to-right direction of \cref{lem:bridge}, using its
four simultaneous obligations: old-to-old reflection, the two path
memberships, and the bridge formula. All path judgments are in $\psi$.
Reflexivity gives the old-to-old obligation and no pair between two distinct
fresh spine variables. For transitivity, an intermediate source forces the
incoming pair to be its identity, and an intermediate sink forces the outgoing
pair to be its identity, by \eqref{eq:g30}. These cases reduce to a subderivation.
An old intermediate variable preserves \eqref{eq:g32} by upward or downward closure.
For an $R_j$-to-$Z_{j'}$ pair through an old $v$, \eqref{eq:g32} gives
$v\in U^R_j\cap L^Z_{j'}$, which is \eqref{eq:g33} with $\alpha=j,\gamma=j'$.

In a decomposition concluding at $R_j$, $j>0$, the lower bounding
literal must be the unique $R$-link in which $R_j$ is a child, and
the selected child position must be $w[j-1]$. For an old right endpoint,
the upper literal is old. The induction hypothesis at the root pair and
the one-step rule \eqref{eq:g22} give membership in $U^R_j$. The old-to-$Z_{j'}$
case is dual. If both endpoints are $R_j,Z_{j'}$, their two unique parent
links force
$w[j-1]=w[j'-1]$. Apply induction to their roots and append this common
letter to the two equal factors in \eqref{eq:g33}. When one parent index is zero, its root is old. If the lower parent is
$R_0$ and the upper parent is $Z_{j'-1}$, \eqref{eq:g32} puts
$R_0$ in $L^Z_{j'-1}$, while reflexivity puts it in $U^R_0$.
This supplies a bridge with empty descent before the final letter.
The other orientation is symmetric. If both parents have index zero,
separation reflects their pair to $D_\psi(R_0,Z_0)$, and $Z_0$ lies
in $U^R_0\cap L^Z_0$.
An old-to-old decomposition has already been dealt with in \eqref{eq:g30}.
If the selected child is off-spine, the conclusion has a bottom filler on
the left or a top filler on the right, so no membership or bridge obligation
applies. The terminal lower bound at $W_m$ has only bottom-filler children;
the terminal upper bound at $Y_m$ has only top-filler children. Thus their
decomposition consequences are covered by the same cases.

This covers every closure rule. Because positive-index $Q$- and
$Y$-variables are sinks, no bridge passes through them, and the fillers
create no further bridges.

It remains to justify zero endpoints in \eqref{eq:g33}. If $j=0$, \eqref{eq:g32} says
$R_0\in L^Z_{j'}$; this gives the formula with $\alpha=0,\gamma=j',v=R_0$.
Conversely, its factor-length equation forces $\gamma=j'$, and a member of
$U^R_0\cap L^Z_{j'}$ yields that same membership by downward closure.
The case $j'=0$ is symmetric. Thus the induction covers every cut. 

\subsection{The Remaining Clash Cases}
We complete the classification in \cref{thm:clashes}. Bound families refer
to literal provenance, including new literals whose cut-zero roots are old.
If the lower bound is old and the upper bound is on $Q_{j'}$, an old top
gives a member of $\mathsf{Top}\cap L^Q_{j'}$. By \eqref{eq:g25} and flip closure, its variable flip is
in $\mathsf{Bot}\cap U_{j'}$, giving bot-readiness. An old constructor
can clash only against $Q_m=\bot$; variable flip and order reversal turn its membership in
$L^Q_m$ into an upper constructor in $U_m$, giving (b).
An old lower bound clashes against a $Y_{j'}$-bound only when it is top;
\eqref{eq:g32} then gives $\operatorname{top}(L_{j'})$.
For a $W_j$-bound against an old upper bound, the latter must be bottom.
Its root belongs to $U^W_j=\flip L_j$; flipping gives an old top in $L_j$,
again top-readiness.

For a $W_j$-bound against a $Q_{j'}$-bound, incompatibility forces $j'=m$.
The bridge gives cuts $\alpha,\gamma$ with
\[
 j-\alpha=m-\gamma,\quad w[\alpha:j]=w[\gamma:m],\quad
 \flip L_\alpha\cap\flip U_\gamma\ne\varnothing.
\]
Flipping a common element gives $L_\alpha\cap U_\gamma\ne\varnothing$.
Since $j\le m$, $\alpha\le\gamma$. Equality gives reflexivity readiness;
otherwise take $s=\alpha,e=\gamma$. Then $j=s+m-e$ and the factor
equality is exactly $w[e:]\preceq w[s:]$, with cross admission.
This includes $e=m$, when descent is empty. A $W$-bound against a
$Y$-bound always has the compatible labels $f,f$, even at terminal cuts.
Its closure bridges may exist, but their decomposition consequences are
already covered by the full invariant in \cref{lem:bridge}.


\subsection{Readiness and Child Events Give Clashes}
Suppose a readiness or child event occurs. If $v\in U_j\cap\mathsf{Bot}$,
\eqref{eq:g32} gives $D(P_j,v)$, and the lower label on $P_j$, either $f$
or $\top$, clashes with bottom. If
$v\in L_j\cap\mathsf{Top}$, then $D(v,Y_j)$ clashes with the upper
constructor on $Y_j$. If $v\in U_j\cap L_j$, then
$D(P_j,Y_j)$; decomposing down the common remaining suffix gives
$D(P_m,Y_m)$, a top/constructor clash. An upper constructor in $U_m$
likewise clashes with $P_m=\top$.
